\documentclass[11pt,a4paper]{article}
\usepackage[T2A]{fontenc}
\usepackage[utf8]{inputenc}
\usepackage[russian]{babel}
\usepackage{geometry}
\usepackage{amsmath,amssymb,mathtools}
\usepackage{array,booktabs,longtable}
\usepackage{xcolor}
\usepackage{hyperref}
\usepackage{tikz}
\geometry{left=21mm,right=18mm,top=20mm,bottom=20mm}
\definecolor{reviewteal}{RGB}{20,105,111}
\hypersetup{colorlinks=true,linkcolor=reviewteal,urlcolor=reviewteal}
\setlength{\parindent}{0pt}
\setlength{\parskip}{5pt}
\setlength{\LTpre}{7pt}
\setlength{\LTpost}{7pt}
\newcommand{\reviewblock}[2]{\par\medskip\noindent\textbf{#1}\par\noindent#2\par}
\newcommand{\errorblock}[1]{\par\medskip\noindent\textcolor{red}{\textbf{Ошибка:}} #1\par}
\newcommand{\ideablock}[1]{\par\medskip\noindent\textcolor{blue!60!black}{\textbf{Ключевая идея:}} #1\par}
\newcommand{\resultblock}[1]{\par\medskip\noindent\textcolor{teal!60!black}{\textbf{Итог:}} #1\par}
\begin{document}
\thispagestyle{empty}
\vspace*{0.28\textheight}
\begin{center}
{\LARGE\bfseries Ревью домашней работы\par}
\vspace{2.3cm}
{\large Автор: Лёвин Артём Александрович\par}
\vspace{0.45cm}
{9 октября 2026 года\par}
\end{center}
\vfill
\begin{center}
\begin{tikzpicture}[x=1cm,y=1cm]
\draw[reviewteal!55,line width=0.55pt] (-2.4,0) .. controls (-1.3,0.5) and (-0.8,-0.5) .. (0,0)
.. controls (0.8,0.5) and (1.3,-0.5) .. (2.4,0);
\end{tikzpicture}
\end{center}
\newpage
\section*{Сводная таблица заданий, требующих разбора}
\small
\begin{longtable}{@{}p{0.065\linewidth}p{0.21\linewidth}p{0.28\linewidth}p{0.145\linewidth}p{0.18\linewidth}@{}}
\toprule
\textbf{№} & \textbf{Тема} & \textbf{Суть ошибки} & \textbf{Ваш ответ} & \textbf{Правильный ответ}\\
\midrule
\endhead
\hyperref[task:5]{5} & Движение по окружности & Решение не приведено & не указан & $5\ \text{м/с}^2$\\
\midrule
\hyperref[task:6]{6} & Закон всемирного тяготения & Решение не приведено & не указан & $1$\\
\midrule
\hyperref[task:7]{7} & Закон сохранения импульса & Не указано направление & 0,25 м/с & $0{,}25\ \text{м/с вправо}$\\
\midrule
\hyperref[task:8]{8} & Закон сохранения энергии & Нет вычисления скорости & не указан & $10\ \text{м/с}$\\
\midrule
\hyperref[task:10]{10} & Закон Архимеда & Решение не приведено & не указан & $75\%$\\
\midrule
\hyperref[task:11]{11} & Пружинный маятник & Решение не приведено & не указан & $2\ \text{раза}$\\
\midrule
\hyperref[task:18]{18} & Влажность воздуха & Решение не приведено & не указан & $50\%$\\
\midrule
\hyperref[task:19]{19} & Закон Кулона & Решение не приведено & не указан & $0{,}9\ \text{мН}$\\
\midrule
\hyperref[task:20]{20} & Напряжённость электрического поля & Решение не приведено & не указан & $400\ \text{Н/Кл}$\\
\midrule
\hyperref[task:21]{21} & Работа электростатического поля & Решение не приведено & не указан & $0{,}3\ \text{мДж}$\\
\midrule
\hyperref[task:22]{22} & Энергия конденсатора & Решение не приведено & не указан & $0{,}1\ \text{Дж}$\\
\midrule
\hyperref[task:24]{24} & Полная цепь постоянного тока & Решение не приведено & не указан & $10\ \text{В}$\\
\midrule
\hyperref[task:26]{26} & Планирование эксперимента & Нет обоснования выбора & 1 и 2 & $1\ \text{и}\ 2$\\
\midrule
\hyperref[task:27]{27} & Сила Ампера (опережающий блок) & Решение не приведено & не указан & $0{,}4\ \text{Н}$\\
\midrule
\hyperref[task:28]{28} & Электромагнитная индукция (опережающий блок) & Решение не приведено & не указан & $0{,}04\ \text{В}$\\
\midrule
\hyperref[task:29]{29} & Тонкая линза (опережающий блок) & Решение не приведено & не указан & $60\ \text{см}$\\
\midrule
\hyperref[task:30]{30} & Энергия фотона (опережающий блок) & Решение не приведено & не указан & $3{,}96$\\
\bottomrule
\end{longtable}
\normalsize
\newpage
\thispagestyle{empty}
\vspace*{0.4\textheight}
\begin{center}
{\Huge\bfseries Разбор ошибок}
\end{center}
\vfill\newpage
\phantomsection\label{task:5}

\section*{Задание № 5}

\reviewblock{Текст задания.}{Автомобиль движется с постоянной по модулю скоростью $15\ \text{м/с}$ по окружности радиусом $45\ \text{м}$. Найдите модуль центростремительного ускорения автомобиля.}

\reviewblock{Ваше решение.}{Под номером 5 есть отметка вопроса, но формулы, вычислений и ответа нет.}

\reviewblock{Ваш ответ.}{Не указан.}

\errorblock{Не найдено центростремительное ускорение.}

\reviewblock{Где именно возникает ошибка.}{Даны скорость и радиус, однако переход к вычислению ускорения отсутствует; первым недостающим шагом является запись закона $a_{\text{ц}}=v^2/R$.}

\reviewblock{Почему этот шаг неверен или неполон.}{Постоянство \emph{модуля} скорости не означает отсутствия ускорения: направление вектора скорости изменяется непрерывно. Ускорение направлено к центру окружности. Для него нужна именно формула $v^2/R$, а не $v/t$ без известного времени.}

\ideablock{Модуль ускорения при равномерном движении по окружности определяется квадратом скорости, делённым на радиус.}

\reviewblock{Исправление от последнего верного шага.}{От исходных данных $v=15\ \text{м/с}$ и $R=45\ \text{м}$ нужно перейти к формуле $a_{\text{ц}}=v^2/R$ и подставить известные значения.}

\reviewblock{Полное правильное решение.}{Скорость по модулю постоянна, поэтому тангенциальная составляющая ускорения равна нулю. Но скорость меняет направление, значит, имеется центростремительное ускорение:
\[a_{\text{ц}}=\frac{v^2}{R}=\frac{15^2}{45}=\frac{225}{45}=5\ \text{м/с}^2.\]
Единица получается из $(\text{м/с})^2/\text{м}=\text{м/с}^2$.}

\reviewblock{Проверка.}{Подстановка в обратное равенство $a_{\text{ц}}R=v^2$ даёт $5\cdot45=225=15^2$.}

\resultblock{Модуль центростремительного ускорения равен $5\ \text{м/с}^2$.}

\reviewblock{Задача на закрепление.}{Автомобиль движется по окружности радиусом $32\ \text{м}$ со скоростью $12\ \text{м/с}$. Найдите модуль центростремительного ускорения.}
\newpage
\phantomsection\label{task:6}

\section*{Задание № 6}

\reviewblock{Текст задания.}{Масса планеты в $9$ раз больше массы Земли, а радиус в $3$ раза больше радиуса Земли. Найдите отношение ускорений свободного падения у поверхностей планеты и Земли.}

\reviewblock{Ваше решение.}{Под номером 6 указана только отметка вопроса; формула и ответ отсутствуют.}

\reviewblock{Ваш ответ.}{Не указан.}

\errorblock{Не установлена зависимость ускорения свободного падения от массы и радиуса планеты.}

\reviewblock{Где именно возникает ошибка.}{Исходные отношения $M_{\text{п}}/M_{\text{З}}=9$ и $R_{\text{п}}/R_{\text{З}}=3$ не использованы; первым недостающим шагом является запись $g=GM/R^2$.}

\reviewblock{Почему этот шаг неверен или неполон.}{Ускорение свободного падения пропорционально массе планеты и обратно пропорционально \emph{квадрату} её радиуса. Увеличение радиуса в три раза уменьшает ускорение при той же массе в девять раз; увеличение массы в девять раз компенсирует этот эффект.}

\ideablock{При сравнении двух планет удобно разделить две формулы $g=GM/R^2$: гравитационная постоянная сократится.}

\reviewblock{Исправление от последнего верного шага.}{Составим дробь $g_{\text{п}}/g_{\text{З}}$ и подставим отношения масс и радиусов, обязательно возведя отношение радиусов в квадрат.}

\reviewblock{Полное правильное решение.}{Для поверхности сферической планеты $g=GM/R^2$. Тогда
\[\frac{g_{\text{п}}}{g_{\text{З}}}=\frac{GM_{\text{п}}/R_{\text{п}}^2}{GM_{\text{З}}/R_{\text{З}}^2}=\frac{M_{\text{п}}}{M_{\text{З}}}\left(\frac{R_{\text{З}}}{R_{\text{п}}}\right)^2.\]
По условию $M_{\text{п}}/M_{\text{З}}=9$ и $R_{\text{п}}/R_{\text{З}}=3$, поэтому
\[\frac{g_{\text{п}}}{g_{\text{З}}}=9\left(\frac13\right)^2=\frac99=1.\]}

\reviewblock{Проверка.}{Множитель массы равен $9$, множитель радиуса равен $1/9$; их произведение равно единице, что согласуется с размерностью искомого отношения.}

\resultblock{Ускорения свободного падения у поверхностей двух планет равны; искомое отношение равно $1$.}

\reviewblock{Задача на закрепление.}{Масса планеты в $16$ раз больше массы Земли, а радиус в $2$ раза больше. Найдите отношение ускорений свободного падения у поверхностей этих планет.}
\newpage
\phantomsection\label{task:7}

\section*{Задание № 7}

\reviewblock{Текст задания.}{Тележка массой $2\ \text{кг}$ движется вправо со скоростью $4\ \text{м/с}$ и сталкивается с тележкой массой $6\ \text{кг}$, движущейся влево со скоростью $1\ \text{м/с}$. После удара тележки сцепляются. Найдите модуль и направление их общей скорости.}

\reviewblock{Ваше решение.}{Записано $p=mv$; далее $2\cdot4-6=8v$, затем $2=8v$ и $v=0{,}25\ \text{м/с}$. Арифметика и найденный модуль скорости верны.}

\reviewblock{Ваш ответ.}{$0{,}25\ \text{м/с}$; направление не записано.}

\errorblock{Ответ неполон: в условии отдельно требуется направление движения.}

\reviewblock{Где именно возникает ошибка.}{Последний верный шаг: $2=8v$, откуда $v=0{,}25\ \text{м/с}$. Первый неполный шаг --- заключительная запись ответа без слов «вправо».}

\reviewblock{Почему этот шаг неверен или неполон.}{Импульс --- векторная величина. При движении по одной прямой его проекции имеют знаки. Положительный результат скорости после решения уравнения указывает на движение в положительном направлении оси. Без определения оси и указания направления численное значение скорости не полностью отвечает на вопрос.}

\ideablock{Выбрать направление вправо положительным и сохранить знаки скоростей до окончательного ответа.}

\reviewblock{Исправление от последнего верного шага.}{Ваша численная часть сохраняется без изменения. Нужно явно записать: $v_1=+4\ \text{м/с}$, $v_2=-1\ \text{м/с}$, поэтому $v=+0{,}25\ \text{м/с}$; знак плюс означает движение вправо.}

\reviewblock{Полное правильное решение.}{Выберем ось $Ox$ вправо. По закону сохранения импульса для неупругого столкновения:
\[m_1v_1+m_2v_2=(m_1+m_2)v.\]
Здесь $m_1=2\ \text{кг}$, $v_1=+4\ \text{м/с}$, $m_2=6\ \text{кг}$, $v_2=-1\ \text{м/с}$. Получаем
\[2\cdot4+6\cdot(-1)=(2+6)v,\]
\[8-6=8v,\qquad v=\frac28=+0{,}25\ \text{м/с}.\]
Знак «плюс» означает движение вправо. Модуль скорости равен $0{,}25\ \text{м/с}$.}

\reviewblock{Проверка.}{Импульс до удара равен $8-6=2\ \text{кг}\cdot\text{м/с}$ вправо; после удара $8\cdot0{,}25=2\ \text{кг}\cdot\text{м/с}$ также вправо.}

\resultblock{После сцепления тележки движутся вправо со скоростью $0{,}25\ \text{м/с}$.}

\reviewblock{Задача на закрепление.}{Тележка массой $3\ \text{кг}$ движется вправо со скоростью $5\ \text{м/с}$ и сцепляется с тележкой массой $2\ \text{кг}$, движущейся влево со скоростью $2\ \text{м/с}$. Найдите модуль и направление общей скорости.}
\newpage
\phantomsection\label{task:8}

\section*{Задание № 8}

\reviewblock{Текст задания.}{Тело начинает скользить без начальной скорости с высоты $5\ \text{м}$ по гладкой неподвижной горке. Потерями механической энергии пренебречь. Найдите скорость тела у основания горки.}

\reviewblock{Ваше решение.}{Есть схематический рисунок наклонной горки, отмечены высота и силы; расчёт скорости и окончательный ответ отсутствуют.}

\reviewblock{Ваш ответ.}{Не указан.}

\errorblock{Решение остановлено на рисунке: закон сохранения механической энергии не применён.}

\reviewblock{Где именно возникает ошибка.}{Последний содержательный шаг --- схема с указанием высоты $h=5\ \text{м}$. Не выполнен переход к уравнению $mgh=mv^2/2$, которое позволяет найти скорость.}

\reviewblock{Почему этот шаг неверен или неполон.}{У гладкой неподвижной горки сила нормальной реакции не совершает работы вдоль перемещения. Трением пренебрегаем. Поэтому потенциальная энергия тела уменьшается ровно на величину прироста его кинетической энергии. Форма горки и масса тела на конечную скорость не влияют.}

\ideablock{Начальную потенциальную энергию относительно основания горки приравнять конечной кинетической энергии.}

\reviewblock{Исправление от последнего верного шага.}{К рисунку следует добавить $v_0=0$, $E_{\text{нач}}=mgh$, $E_{\text{кон}}=mv^2/2$, затем приравнять энергии и сократить массу.}

\reviewblock{Полное правильное решение.}{Выберем нулевой уровень потенциальной энергии у основания горки. В начале движения тело покоится на высоте $h$: $E_{\text{нач}}=mgh+0$. У основания высота равна нулю: $E_{\text{кон}}=0+mv^2/2$. Так как потерь нет,
\[mgh=\frac{mv^2}{2}.\]
Масса отлична от нуля, сократим её:
\[gh=\frac{v^2}{2},\qquad v=\sqrt{2gh}.\]
Подставляем $g=10\ \text{м/с}^2$, $h=5\ \text{м}$:
\[v=\sqrt{2\cdot10\cdot5}=\sqrt{100}=10\ \text{м/с}.\]}

\reviewblock{Проверка.}{Для полученной скорости $v^2/2=100/2=50\ \text{Дж/кг}$, а $gh=10\cdot5=50\ \text{Дж/кг}$; удельные энергии равны.}

\resultblock{Скорость у основания горки составляет $10\ \text{м/с}$.}

\reviewblock{Задача на закрепление.}{Тело из состояния покоя скользит по гладкой горке с высоты $7{,}2\ \text{м}$. Найдите скорость тела у основания, пренебрегая потерями энергии.}
\newpage
\phantomsection\label{task:10}

\section*{Задание № 10}

\reviewblock{Текст задания.}{Однородный брусок плотностью $750\ \text{кг/м}^3$ плавает в воде плотностью $1000\ \text{кг/м}^3$. Какая доля объёма бруска находится под водой? Ответ дайте в процентах.}

\reviewblock{Ваше решение.}{Под номером 10 есть отметка вопроса; вычислений нет.}

\reviewblock{Ваш ответ.}{Не указан.}

\errorblock{Не определена доля погружённого объёма.}

\reviewblock{Где именно возникает ошибка.}{В исходных данных известны плотности, однако не записано условие плавания $F_{\text{А}}=mg$. Следовательно, объём погружённой части не найден.}

\reviewblock{Почему этот шаг неверен или неполон.}{Для плавающего покоящегося тела архимедова сила точно уравновешивает тяжесть. Архимедова сила зависит от объёма вытесненной воды, то есть погружённой части бруска, а сила тяжести --- от полного объёма бруска.}

\ideablock{Сравнить силы и сократить одинаковые величины $g$ и полный объём, получив отношение плотностей.}

\reviewblock{Исправление от последнего верного шага.}{Из $\rho_{\text{в}}gV_{\text{погр}}=\rho_{\text{б}}gV$ получить $V_{\text{погр}}/V=\rho_{\text{б}}/\rho_{\text{в}}$.}

\reviewblock{Полное правильное решение.}{На плавающий брусок действуют сила тяжести $mg$ вниз и архимедова сила $F_{\text{А}}$ вверх. Равновесие требует $F_{\text{А}}=mg$. По закону Архимеда
\[F_{\text{А}}=\rho_{\text{в}}gV_{\text{погр}},\]
а масса всего бруска $m=\rho_{\text{б}}V$. Значит,
\[\rho_{\text{в}}gV_{\text{погр}}=\rho_{\text{б}}gV.\]
Сократим $g$ и выразим искомую долю:
\[\frac{V_{\text{погр}}}{V}=\frac{\rho_{\text{б}}}{\rho_{\text{в}}}=\frac{750}{1000}=0{,}75.\]
Переведём долю в проценты: $0{,}75\cdot100\%=75\%$.}

\reviewblock{Проверка.}{Доля $75\%$ меньше $100\%$, что правдоподобно: плотность бруска меньше плотности воды, поэтому он действительно плавает.}

\resultblock{Под водой находится $75\%$ объёма бруска.}

\reviewblock{Задача на закрепление.}{Деревянный брусок плотностью $600\ \text{кг/м}^3$ плавает в жидкости плотностью $800\ \text{кг/м}^3$. Найдите долю погружённого объёма в процентах.}
\newpage
\phantomsection\label{task:11}

\section*{Задание № 11}

\reviewblock{Текст задания.}{Массу груза на пружине увеличили в $4$ раза, жёсткость пружины не изменилась. Во сколько раз изменился период свободных колебаний?}

\reviewblock{Ваше решение.}{Под номером 11 оставлена отметка вопроса; решение отсутствует.}

\reviewblock{Ваш ответ.}{Не указан.}

\errorblock{Не использована зависимость периода от массы груза.}

\reviewblock{Где именно возникает ошибка.}{Дано, что $m_2=4m_1$ и $k_2=k_1$, но не записана формула $T=2\pi\sqrt{m/k}$, позволяющая сравнить периоды.}

\reviewblock{Почему этот шаг неверен или неполон.}{Период пружинного маятника зависит от \emph{квадратного корня} из массы. Поэтому увеличение массы в четыре раза не приводит к четырёхкратному увеличению периода.}

\ideablock{Записать отношение $T_2/T_1$ и сократить постоянные множители и жёсткость.}

\reviewblock{Исправление от последнего верного шага.}{Пользуясь $T=2\pi\sqrt{m/k}$, выразить $T_2/T_1=\sqrt{m_2/m_1}$.}

\reviewblock{Полное правильное решение.}{Период колебаний груза на пружине при малых колебаниях описывается формулой
\[T=2\pi\sqrt{\frac{m}{k}}.\]
Начальный и конечный периоды:
\[T_1=2\pi\sqrt{\frac{m_1}{k}},\qquad T_2=2\pi\sqrt{\frac{m_2}{k}}.\]
Разделим второе выражение на первое:
\[\frac{T_2}{T_1}=\sqrt{\frac{m_2}{m_1}}=\sqrt4=2.\]
Следовательно, $T_2=2T_1$: период увеличится в два раза.}

\reviewblock{Проверка.}{Если массу заменить на $4m$, то $2\pi\sqrt{4m/k}=2\cdot2\pi\sqrt{m/k}$, как и найдено.}

\resultblock{Период увеличится в $2$ раза.}

\reviewblock{Задача на закрепление.}{Массу груза пружинного маятника увеличили в $9$ раз, не меняя жёсткость пружины. Во сколько раз изменится период колебаний?}
\newpage
\phantomsection\label{task:18}

\section*{Задание № 18}

\reviewblock{Текст задания.}{Плотность насыщенного водяного пара при данной температуре равна $17{,}3\ \text{г/м}^3$, а фактическая плотность водяного пара в воздухе --- $8{,}65\ \text{г/м}^3$. Найдите относительную влажность в процентах.}

\reviewblock{Ваше решение.}{Под номером 18 поставлен вопросительный знак; численного решения нет.}

\reviewblock{Ваш ответ.}{Не указан.}

\errorblock{Не найдено отношение фактической плотности пара к плотности насыщенного пара.}

\reviewblock{Где именно возникает ошибка.}{Две плотности указаны в условии, но формула для относительной влажности $\varphi=\rho/\rho_{\text{нас}}\cdot100\%$ не применена.}

\reviewblock{Почему этот шаг неверен или неполон.}{Относительная влажность показывает, какую долю от максимально возможного при данной температуре содержания водяного пара составляет фактическое содержание. Плотности должны относиться к одной и той же температуре.}

\ideablock{Разделить фактическую плотность на плотность насыщенного пара и перевести полученное отношение в проценты.}

\reviewblock{Исправление от последнего верного шага.}{Взять $\rho=8{,}65\ \text{г/м}^3$, $\rho_{\text{нас}}=17{,}3\ \text{г/м}^3$; одинаковые единицы при делении сократятся.}

\reviewblock{Полное правильное решение.}{Определение относительной влажности при фиксированной температуре:
\[\varphi=\frac{\rho}{\rho_{\text{нас}}}\cdot100\%.\]
Подставим исходные данные:
\[\varphi=\frac{8{,}65}{17{,}3}\cdot100\%=0{,}5\cdot100\%=50\%.\]
Фактическая плотность ровно в два раза меньше плотности насыщенного пара, значит, влажность равна половине от $100\%$.}

\reviewblock{Проверка.}{Проверка: $17{,}3\cdot0{,}5=8{,}65$. Влажность лежит между $0\%$ и $100\%$.}

\resultblock{Относительная влажность воздуха составляет $50\%$.}

\reviewblock{Задача на закрепление.}{Фактическая плотность водяного пара равна $9\ \text{г/м}^3$, а плотность насыщенного пара при той же температуре --- $15\ \text{г/м}^3$. Найдите относительную влажность.}
\newpage
\phantomsection\label{task:19}

\section*{Задание № 19}

\reviewblock{Текст задания.}{Заряды $2\cdot10^{-8}\ \text{Кл}$ и $5\cdot10^{-8}\ \text{Кл}$ находятся в вакууме на расстоянии $0{,}10\ \text{м}$. Найдите модуль силы взаимодействия в миллиньютонах.}

\reviewblock{Ваше решение.}{Под номером 19 указана отметка вопроса; формула и ответ отсутствуют.}

\reviewblock{Ваш ответ.}{Не указан.}

\errorblock{Не применён закон Кулона и не выполнен перевод ньютонов в миллиньютоны.}

\reviewblock{Где именно возникает ошибка.}{Известны оба заряда и расстояние, но вычисление не начинается с $F=k|q_1q_2|/r^2$.}

\reviewblock{Почему этот шаг неверен или неполон.}{Сила электрического взаимодействия точечных зарядов в вакууме прямо пропорциональна произведению модулей зарядов и обратно пропорциональна квадрату расстояния между ними. Для модуля силы знак зарядов значения не имеет.}

\ideablock{Подставить заряды в кулонах и расстояние в метрах; возвести расстояние в квадрат; затем умножить результат в ньютонах на $1000$ для перехода к миллиньютонам.}

\reviewblock{Исправление от последнего верного шага.}{Используем $k=9\cdot10^9\ \text{Н}\cdot\text{м}^2/\text{Кл}^2$ и $r^2=(0{,}10)^2=0{,}01\ \text{м}^2$.}

\reviewblock{Полное правильное решение.}{По закону Кулона
\[F=k\frac{|q_1q_2|}{r^2}.\]
Произведение модулей зарядов:
\[|q_1q_2|=(2\cdot10^{-8})(5\cdot10^{-8})=10\cdot10^{-16}=10^{-15}\ \text{Кл}^2.\]
Квадрат расстояния равен $0{,}10^2=0{,}01\ \text{м}^2$. Следовательно,
\[F=9\cdot10^9\cdot\frac{10^{-15}}{0{,}01}=9\cdot10^{-4}\ \text{Н}.\]
Так как $1\ \text{Н}=1000\ \text{мН}$, получаем
\[F=9\cdot10^{-4}\cdot10^3=0{,}9\ \text{мН}.\]}

\reviewblock{Проверка.}{Проверка порядка величины: заряды порядка $10^{-8}\ \text{Кл}$ на расстоянии порядка $10^{-1}\ \text{м}$ дают силу порядка $10^{-3}\ \text{Н}$, что соответствует $0{,}9\ \text{мН}$.}

\resultblock{Модуль силы взаимодействия равен $0{,}9\ \text{мН}$.}

\reviewblock{Задача на закрепление.}{В вакууме находятся заряды $3\cdot10^{-8}\ \text{Кл}$ и $4\cdot10^{-8}\ \text{Кл}$ на расстоянии $0{,}20\ \text{м}$. Найдите модуль силы их взаимодействия в миллиньютонах.}
\newpage
\phantomsection\label{task:20}

\section*{Задание № 20}

\reviewblock{Текст задания.}{Найдите модуль напряжённости электрического поля точечного заряда $4\ \text{нКл}$ на расстоянии $0{,}30\ \text{м}$ от него.}

\reviewblock{Ваше решение.}{Под номером 20 решения и ответа нет.}

\reviewblock{Ваш ответ.}{Не указан.}

\errorblock{Не применена формула напряжённости поля точечного заряда.}

\reviewblock{Где именно возникает ошибка.}{Величина заряда и расстояние заданы, но не сделан переход $E=k|q|/r^2$; не выполнен перевод нанокулонов в кулоны.}

\reviewblock{Почему этот шаг неверен или неполон.}{Напряжённость поля точечного заряда рассчитывается по закону, обратному квадрату расстояния. При переходе к СИ $1\ \text{нКл}=10^{-9}\ \text{Кл}$.}

\ideablock{Перевести заряд в кулоны, найти квадрат расстояния и вычислить $E=k|q|/r^2$.}

\reviewblock{Исправление от последнего верного шага.}{Используем $q=4\cdot10^{-9}\ \text{Кл}$ и $r^2=0{,}09\ \text{м}^2$.}

\reviewblock{Полное правильное решение.}{Формула для модуля напряжённости:
\[E=k\frac{|q|}{r^2}.\]
Из условия: $|q|=4\ \text{нКл}=4\cdot10^{-9}\ \text{Кл}$, $r=0{,}30\ \text{м}$. Подставляем:
\[E=\frac{9\cdot10^9\cdot4\cdot10^{-9}}{(0{,}30)^2}=\frac{36}{0{,}09}=400\ \text{Н/Кл}.\]
Для положительного заряда вектор поля был бы направлен от заряда, но вопрос требует только модуль.}

\reviewblock{Проверка.}{Умножение $E r^2$ даёт $400\cdot0{,}09=36=k|q|$ в соответствующих единицах.}

\resultblock{Модуль напряжённости составляет $400\ \text{Н/Кл}$.}

\reviewblock{Задача на закрепление.}{Найдите модуль напряжённости поля заряда $3\ \text{нКл}$ на расстоянии $0{,}30\ \text{м}$ в вакууме.}
\newpage
\phantomsection\label{task:21}

\section*{Задание № 21}

\reviewblock{Текст задания.}{Положительный заряд $3\ \text{мкКл}$ перемещается из точки с потенциалом $120\ \text{В}$ в точку с потенциалом $20\ \text{В}$. Найдите работу электрического поля в миллиджоулях.}

\reviewblock{Ваше решение.}{Под номером 21 формулы и ответа нет.}

\reviewblock{Ваш ответ.}{Не указан.}

\errorblock{Не вычислена разность потенциалов и работа поля.}

\reviewblock{Где именно возникает ошибка.}{Из данных $\varphi_1=120\ \text{В}$ и $\varphi_2=20\ \text{В}$ не получено выражение $A=q(\varphi_1-\varphi_2)$.}

\reviewblock{Почему этот шаг неверен или неполон.}{Работа электростатического поля при перемещении заряда между точками зависит только от их потенциалов. Для положительного заряда при переходе к меньшему потенциалу работа поля положительна.}

\ideablock{Найти разность потенциалов, перевести микрокулоны в кулоны и миллиджоули в джоули при вычислениях.}

\reviewblock{Исправление от последнего верного шага.}{Здесь $\Delta\varphi=120-20=100\ \text{В}$, $q=3\cdot10^{-6}\ \text{Кл}$.}

\reviewblock{Полное правильное решение.}{Работа электростатического поля равна
\[A=q(\varphi_1-\varphi_2).\]
Переведём заряд: $3\ \text{мкКл}=3\cdot10^{-6}\ \text{Кл}$. Вычислим разность потенциалов:
\[\varphi_1-\varphi_2=120-20=100\ \text{В}.\]
Тогда
\[A=3\cdot10^{-6}\cdot100=3\cdot10^{-4}\ \text{Дж}.\]
Поскольку $1\ \text{мДж}=10^{-3}\ \text{Дж}$,
\[A=\frac{3\cdot10^{-4}}{10^{-3}}=0{,}3\ \text{мДж}.\]}

\reviewblock{Проверка.}{Знак положителен, поскольку положительный заряд перемещён к меньшему потенциалу. Проверка единиц: $\text{Кл}\cdot\text{В}=\text{Дж}$.}

\resultblock{Работа поля равна $+0{,}3\ \text{мДж}$.}

\reviewblock{Задача на закрепление.}{Положительный заряд $5\ \text{мкКл}$ переместился из точки с потенциалом $90\ \text{В}$ в точку с потенциалом $30\ \text{В}$. Найдите работу электрического поля в миллиджоулях.}
\newpage
\phantomsection\label{task:22}

\section*{Задание № 22}

\reviewblock{Текст задания.}{Конденсатор ёмкостью $20\ \text{мкФ}$ заряжен до напряжения $100\ \text{В}$. Найдите энергию его электрического поля.}

\reviewblock{Ваше решение.}{Под номером 22 нет решения и ответа.}

\reviewblock{Ваш ответ.}{Не указан.}

\errorblock{Не применена формула энергии заряженного конденсатора.}

\reviewblock{Где именно возникает ошибка.}{Даны ёмкость и напряжение, но отсутствует первый расчётный шаг $W=CU^2/2$.}

\reviewblock{Почему этот шаг неверен или неполон.}{Энергия заряженного конденсатора пропорциональна ёмкости и \emph{квадрату} напряжения. Множитель $1/2$ обязателен. Микрофарады следует перевести в фарады.}

\ideablock{Записать энергию через ёмкость и напряжение, затем перейти к единицам СИ.}

\reviewblock{Исправление от последнего верного шага.}{Подставить $C=20\cdot10^{-6}\ \text{Ф}$ и $U=100\ \text{В}$.}

\reviewblock{Полное правильное решение.}{Используем формулу энергии электрического поля конденсатора:
\[W=\frac{CU^2}{2}.\]
Перевод единиц: $20\ \text{мкФ}=20\cdot10^{-6}\ \text{Ф}$. Квадрат напряжения: $100^2=10^4\ \text{В}^2$. Тогда
\[W=\frac{20\cdot10^{-6}\cdot10^4}{2}=\frac{20\cdot10^{-2}}2=0{,}1\ \text{Дж}.\]
Если забыть деление на два, получится вдвое завышенное значение, поэтому коэффициент необходимо сохранить до конца.}

\reviewblock{Проверка.}{Обратная проверка: $2W/U^2=0{,}2/10000=2\cdot10^{-5}\ \text{Ф}=20\ \text{мкФ}$.}

\resultblock{Энергия электрического поля равна $0{,}1\ \text{Дж}$.}

\reviewblock{Задача на закрепление.}{Конденсатор ёмкостью $50\ \text{мкФ}$ заряжен до напряжения $40\ \text{В}$. Найдите запасённую энергию.}
\newpage
\phantomsection\label{task:24}

\section*{Задание № 24}

\reviewblock{Текст задания.}{Источник имеет ЭДС $12\ \text{В}$ и внутреннее сопротивление $1\ \text{Ом}$. К нему подключён резистор сопротивлением $5\ \text{Ом}$. Найдите напряжение на внешнем резисторе.}

\reviewblock{Ваше решение.}{Под номером 24 оставлена отметка вопроса; расчёта нет.}

\reviewblock{Ваш ответ.}{Не указан.}

\errorblock{Не учтено внутреннее сопротивление источника при расчёте напряжения на нагрузке.}

\reviewblock{Где именно возникает ошибка.}{Из известных $\mathcal{E}$, $r$ и $R$ не найден ток по закону Ома для полной цепи; без этого нельзя надёжно определить напряжение $U=IR$.}

\reviewblock{Почему этот шаг неверен или неполон.}{ЭДС --- это не всегда напряжение на внешнем резисторе. При протекании тока часть ЭДС приходится на падение напряжения на внутреннем сопротивлении источника: $\mathcal{E}=IR+Ir$.}

\ideablock{Найти полный ток $I=\mathcal{E}/(R+r)$, затем вычислить напряжение на нагрузке $U=IR$.}

\reviewblock{Исправление от последнего верного шага.}{Заменить $R+r$ суммой $5+1=6\ \text{Ом}$ и получить ток $I=12/6=2\ \text{А}$.}

\reviewblock{Полное правильное решение.}{По закону Ома для полной цепи
\[I=\frac{\mathcal{E}}{R+r}.\]
Подставим значения:
\[I=\frac{12}{5+1}=\frac{12}{6}=2\ \text{А}.\]
Напряжение на внешнем резисторе:
\[U=IR=2\cdot5=10\ \text{В}.\]
Можно прийти к тому же результату, вычитая внутреннее падение напряжения из ЭДС:
\[U=\mathcal{E}-Ir=12-2\cdot1=10\ \text{В}.\]}

\reviewblock{Проверка.}{Проверка баланса напряжений: $10\ \text{В}$ на резисторе плюс $2\ \text{В}$ внутри источника равняются ЭДС $12\ \text{В}$.}

\resultblock{Напряжение на внешнем резисторе равно $10\ \text{В}$.}

\reviewblock{Задача на закрепление.}{Источник с ЭДС $15\ \text{В}$ и внутренним сопротивлением $2\ \text{Ом}$ подключён к резистору $3\ \text{Ом}$. Найдите напряжение на внешнем резисторе.}
\newpage
\phantomsection\label{task:26}

\section*{Задание № 26}

\reviewblock{Текст задания.}{Для проверки зависимости периода малых колебаний маятника от длины нужно сравнить два набора: 1) $40\ \text{см}$, $100\ \text{г}$, $5^\circ$; 2) $80\ \text{см}$, $100\ \text{г}$, $5^\circ$; 3) $80\ \text{см}$, $200\ \text{г}$, $5^\circ$; 4) $40\ \text{см}$, $100\ \text{г}$, $15^\circ$.}

\reviewblock{Ваше решение.}{Записан только ответ «1 и 2» без объяснения выбора. Сами номера подобраны правильно, но инструкция требует обосновать каждое решение.}

\reviewblock{Ваш ответ.}{«1 и 2»; числовой выбор правильный, обоснование не приведено.}

\errorblock{Не показано, почему сравнивать следует именно наборы 1 и 2.}

\reviewblock{Где именно возникает ошибка.}{Последний правильный шаг --- выбор набора 1 и набора 2. Первый недостающий шаг --- проверка того, что меняется только длина нити, а масса и угол сохраняются.}

\reviewblock{Почему этот шаг неверен или неполон.}{В контролируемом эксперименте для установления влияния одной величины остальные существенные условия оставляют одинаковыми. При сравнении наборов 2 и 3 изменяется масса, а не длина. При сравнении 1 и 4 меняется начальный угол. Только пара 1 и 2 позволяет изолировать влияние длины при одинаковых массе и угле.}

\ideablock{Сравнивать две установки, различающиеся только исследуемым параметром --- длиной нити.}

\reviewblock{Исправление от последнего верного шага.}{Ваш выбор «1 и 2» сохраняется. Добавить пояснение: длины $40$ и $80\ \text{см}$ различаются, но массы в обоих опытах равны $100\ \text{г}$, а углы --- $5^\circ$.}

\reviewblock{Полное правильное решение.}{Исследуемая величина --- период $T$, изменяемый фактор --- длина нити $l$. Для математического маятника при малых колебаниях
\[T=2\pi\sqrt{\frac{l}{g}}.\]
Чтобы установить влияние $l$, нужно сравнить наборы с разными $l$, но одинаковыми остальными условиями. В наборе 1 имеем $l_1=40\ \text{см}$, $m_1=100\ \text{г}$, $\alpha_1=5^\circ$. В наборе 2: $l_2=80\ \text{см}$, $m_2=100\ \text{г}$, $\alpha_2=5^\circ$. Следовательно, изменяется только длина нити, и сравнение корректно. Выбираем наборы $1$ и $2$.}

\reviewblock{Проверка.}{Проверяем по таблице: $40\ne80$, $100=100$, $5^\circ=5^\circ$. Ровно одна характеристика установок меняется.}

\resultblock{Нужные наборы --- $1$ и $2$; теперь выбор обоснован.}

\reviewblock{Задача на закрепление.}{В опыте с маятником имеются наборы: A) $50\ \text{см}$, $150\ \text{г}$, $4^\circ$; Б) $70\ \text{см}$, $150\ \text{г}$, $4^\circ$; В) $70\ \text{см}$, $250\ \text{г}$, $4^\circ$. Какие два набора подходят для проверки зависимости периода от длины?}
\newpage
\phantomsection\label{task:27}

\section*{Задание № 27}

\reviewblock{Текст задания.}{Прямой проводник длиной $0{,}50\ \text{м}$ перпендикулярен линиям однородного магнитного поля с индукцией $0{,}40\ \text{Тл}$. Сила тока равна $2{,}0\ \text{А}$. Найдите модуль силы Ампера.}

\reviewblock{Ваше решение.}{Выписаны исходные данные $l=0{,}5\ \text{м}$, $B=0{,}4\ \text{Тл}$, $I=2\ \text{А}$, но расчёта и ответа нет.}

\reviewblock{Ваш ответ.}{Не указан.}

\errorblock{Не завершён расчёт силы Ампера.}

\reviewblock{Где именно возникает ошибка.}{Последний правильный шаг --- запись исходных данных. Первый отсутствующий шаг --- $F_{\text{А}}=BIl\sin\alpha$.}

\reviewblock{Почему этот шаг неверен или неполон.}{Для прямого проводника модуль магнитной силы зависит от синуса угла между направлением тока и вектором индукции. По условию проводник перпендикулярен линиям поля, поэтому $\alpha=90^\circ$ и $\sin\alpha=1$.}

\ideablock{При перпендикулярности проводника магнитному полю формула упрощается до $F_{\text{А}}=BIl$.}

\reviewblock{Исправление от последнего верного шага.}{Из записанных данных сразу подставить $B=0{,}40$, $I=2{,}0$, $l=0{,}50$ в произведение $BIl$.}

\reviewblock{Полное правильное решение.}{Общая формула силы Ампера:
\[F_{\text{А}}=BIl\sin\alpha.\]
Угол между током и магнитным полем равен $90^\circ$, поэтому $\sin90^\circ=1$. Тогда
\[F_{\text{А}}=0{,}40\cdot2{,}0\cdot0{,}50=0{,}40\ \text{Н}.\]
В задаче требуется модуль силы; её направление по правилу левой руки определять не нужно.}

\reviewblock{Проверка.}{Сила не может превышать $BIl=0{,}4\ \text{Н}$; при угле $90^\circ$ достигается именно это значение.}

\resultblock{Модуль силы Ампера равен $0{,}4\ \text{Н}$.}

\reviewblock{Задача на закрепление.}{Проводник длиной $0{,}30\ \text{м}$ расположен перпендикулярно однородному полю с индукцией $0{,}50\ \text{Тл}$. Ток равен $4\ \text{А}$. Найдите силу Ампера.}
\newpage
\phantomsection\label{task:28}

\section*{Задание № 28}

\reviewblock{Текст задания.}{Площадь одного витка $0{,}020\ \text{м}^2$. Перпендикулярное его плоскости магнитное поле за $0{,}20\ \text{с}$ изменилось от $0{,}10$ до $0{,}50\ \text{Тл}$. Найдите модуль ЭДС индукции.}

\reviewblock{Ваше решение.}{Для задания 28 решения нет.}

\reviewblock{Ваш ответ.}{Не указан.}

\errorblock{Не рассчитано изменение магнитного потока и ЭДС индукции.}

\reviewblock{Где именно возникает ошибка.}{Даны площадь, интервал времени и два значения индукции, но не записан закон Фарадея для модуля ЭДС: $|\mathcal{E}_{\text{и}}|=|\Delta\Phi|/\Delta t$.}

\reviewblock{Почему этот шаг неверен или неполон.}{Магнитный поток через виток равен $\Phi=BS\cos\alpha$, где $\alpha$ --- угол между индукцией и нормалью к плоскости витка. Поле перпендикулярно плоскости, следовательно, направлено вдоль нормали, $|\cos\alpha|=1$ и изменение потока по модулю равно $S|\Delta B|$.}

\ideablock{Сначала найти $|\Delta B|=0{,}50-0{,}10$, затем использовать $|\mathcal{E}_{\text{и}}|=S|\Delta B|/\Delta t$.}

\reviewblock{Исправление от последнего верного шага.}{При постоянной площади и перпендикулярном поле $|\Delta\Phi|=0{,}020\cdot0{,}40\ \text{Вб}$.}

\reviewblock{Полное правильное решение.}{Изменение магнитной индукции:
\[|\Delta B|=|0{,}50-0{,}10|=0{,}40\ \text{Тл}.\]
При указанной ориентации магнитный поток $\Phi=BS$, поэтому
\[|\Delta\Phi|=S|\Delta B|=0{,}020\cdot0{,}40=0{,}008\ \text{Вб}.\]
По закону электромагнитной индукции для одного витка
\[|\mathcal{E}_{\text{и}}|=\frac{|\Delta\Phi|}{\Delta t}=\frac{0{,}008}{0{,}20}=0{,}04\ \text{В}.\]
Знак ЭДС определяет направление индукционного тока, но в вопросе требуется только модуль.}

\reviewblock{Проверка.}{Умножая $0{,}04\ \text{В}$ на $0{,}20\ \text{с}$, получаем изменение потока $0{,}008\ \text{Вб}$.}

\resultblock{Модуль ЭДС индукции равен $0{,}04\ \text{В}$.}

\reviewblock{Задача на закрепление.}{Виток площадью $0{,}030\ \text{м}^2$ расположен перпендикулярно магнитному полю. Индукция за $0{,}15\ \text{с}$ изменилась от $0{,}20$ до $0{,}50\ \text{Тл}$. Найдите модуль ЭДС индукции.}
\newpage
\phantomsection\label{task:29}

\section*{Задание № 29}

\reviewblock{Текст задания.}{Предмет находится на расстоянии $30\ \text{см}$ от собирающей линзы с фокусным расстоянием $20\ \text{см}$. На каком расстоянии от линзы получится действительное изображение? Ответ дайте в сантиметрах.}

\reviewblock{Ваше решение.}{Для задания 29 решения нет.}

\reviewblock{Ваш ответ.}{Не указан.}

\errorblock{Не применена формула тонкой линзы.}

\reviewblock{Где именно возникает ошибка.}{Исходные расстояние до предмета $a=30\ \text{см}$ и фокусное расстояние $f=20\ \text{см}$ не связаны уравнением $1/f=1/a+1/b$.}

\reviewblock{Почему этот шаг неверен или неполон.}{Для собирающей линзы с действительным изображением расстояние до изображения $b$ положительно. Предмет расположен дальше фокуса: $30>20$, значит, действительное изображение существует. Нельзя складывать расстояния $a$ и $f$ напрямую.}

\ideablock{Из формулы линзы выразить $1/b=1/f-1/a$, затем вычислить обратную величину.}

\reviewblock{Исправление от последнего верного шага.}{Использовать одни единицы --- сантиметры; при вычислении общего знаменателя $60$ удобно вычесть дроби.}

\reviewblock{Полное правильное решение.}{Формула тонкой линзы:
\[\frac1f=\frac1a+\frac1b.\]
Ищем расстояние до изображения $b$:
\[\frac1b=\frac1f-\frac1a=\frac1{20}-\frac1{30}=\frac3{60}-\frac2{60}=\frac1{60}\ \text{см}^{-1}.\]
Следовательно,
\[b=60\ \text{см}.\]
Поскольку $a=30\ \text{см}>f=20\ \text{см}$, действительное изображение действительно получается на другой стороне линзы.}

\reviewblock{Проверка.}{Подстановка обратно: $1/30+1/60=2/60+1/60=3/60=1/20$, формула выполняется.}

\resultblock{Действительное изображение расположено в $60\ \text{см}$ от линзы.}

\reviewblock{Задача на закрепление.}{Предмет находится в $45\ \text{см}$ от собирающей линзы с фокусным расстоянием $15\ \text{см}$. Найдите расстояние от линзы до действительного изображения.}
\newpage
\phantomsection\label{task:30}

\section*{Задание № 30}

\reviewblock{Текст задания.}{Частота фотона равна $6{,}0\cdot10^{14}\ \text{Гц}$. Найдите энергию фотона; ответ запишите в единицах $10^{-19}\ \text{Дж}$. Используйте $h=6{,}6\cdot10^{-34}\ \text{Дж}\cdot\text{с}$.}

\reviewblock{Ваше решение.}{Для задания 30 решение и ответ отсутствуют.}

\reviewblock{Ваш ответ.}{Не указан.}

\errorblock{Не применена формула энергии фотона и не выполнено представление результата в требуемых единицах.}

\reviewblock{Где именно возникает ошибка.}{Частота дана, но не записана формула $E=h\nu$. После вычисления энергии необходимо дополнительно выразить её числом единиц $10^{-19}\ \text{Дж}$.}

\reviewblock{Почему этот шаг неверен или неполон.}{Энергия одного фотона пропорциональна частоте электромагнитного излучения. При умножении степеней десяти показатели складываются: $10^{-34}\cdot10^{14}=10^{-20}$.}

\ideablock{Вычислить $E=h\nu$ в джоулях, затем разделить его на $10^{-19}\ \text{Дж}$, чтобы получить число для ответа.}

\reviewblock{Исправление от последнего верного шага.}{Подставить $h=6{,}6\cdot10^{-34}$, $\nu=6{,}0\cdot10^{14}$ и аккуратно перемножить мантиссы и степени.}

\reviewblock{Полное правильное решение.}{Для фотона справедлива формула Планка
\[E=h\nu.\]
Подставляем числа:
\[E=(6{,}6\cdot10^{-34})(6{,}0\cdot10^{14})\ \text{Дж}=39{,}6\cdot10^{-20}\ \text{Дж}.\]
Приведём число к стандартному виду:
\[E=3{,}96\cdot10^{-19}\ \text{Дж}.\]
Условие просит записать ответ \emph{в единицах} $10^{-19}\ \text{Дж}$. Значит, окончательное число --- $3{,}96$.}

\reviewblock{Проверка.}{Проверка размерности: $(\text{Дж}\cdot\text{с})\cdot(1/\text{с})=\text{Дж}$. Обратное вычисление $3{,}96\cdot10^{-19}/(6{,}6\cdot10^{-34})=6{,}0\cdot10^{14}\ \text{Гц}$.}

\resultblock{Энергия фотона равна $3{,}96\cdot10^{-19}\ \text{Дж}$; в указанных единицах ответ: $3{,}96$.}

\reviewblock{Задача на закрепление.}{Найдите энергию фотона с частотой $5{,}0\cdot10^{14}\ \text{Гц}$. Используйте $h=6{,}6\cdot10^{-34}\ \text{Дж}\cdot\text{с}$. Ответ выразите числом единиц $10^{-19}\ \text{Дж}$.}
\par\vfill
\begin{center}
\small Лёвин Артём Александрович эксклюзивно для Ивана Петраченкова
\end{center}
\end{document}